% ---------------------------------------------------------------------
%  Chapter 1 : Overview and Architecture
% ---------------------------------------------------------------------
\chapter{Overview and Architecture}
\label{ch:overview}

\section{What the engine computes}

You hand \Daedalus{} a stochastic differential equation --- an ordinary one in
time, or a partial one in space and time --- of the general semilinear
(Langevin) form
\begin{equation}
  \Dt \phi_a(\mathbf x,t)
   \;=\; -\,\mu_a\,\phi_a + D_a \Lap \phi_a
          \;+\; V_a\big(\{\phi\}\big) \;+\; \xi_a(\mathbf x,t),
  \label{eq:langevin}
\end{equation}
where $\phi_a$ are the fields, $V_a$ collects the (poly)nonlinear interactions,
and $\xi_a$ is a Gaussian noise with a prescribed covariance (white, or
exponentially correlated in time).  In return the engine computes the connected
\term{cumulants} of the fields,
\begin{equation}
  \kappa_k(t_1,\dots,t_k)
    \;=\; \big\langle \phi(t_1)\cdots\phi(t_k)\big\rangle_{\mathrm c},
\end{equation}
to a chosen order in the loop expansion, as exact symbolic-then-numeric
expressions --- with \emph{no} model-specific code.  The same machinery handles
an Ornstein--Uhlenbeck particle, a multi-population nonlinear Hawkes network, a
reaction--diffusion field, and the KPZ equation, because each is just a different
instance of \eqref{eq:langevin}.

\begin{defn}
\term{Model-independence} is the central design principle: every algorithm in the
pipeline reads the structure of the action symbolically and does the right thing,
rather than special-casing a particular equation.  Adding a new model means
writing a short text description (Chapter on theory files); it never means
touching the integrator.
\end{defn}

\section{The MSR--JD formalism in one page}

The engine never works with the SDE \eqref{eq:langevin} directly.  It works with
the equivalent field theory obtained by the
Martin--Siggia--Rose--Janssen--De~Dominicis (\msrjd) construction, which trades a
stochastic equation for a deterministic path integral over \emph{two} fields per
degree of freedom.

For each physical field $\phi_a$ one introduces a \term{response field}
$\tilde\phi_a$ (the engine writes it \code{xt} for a physical field \code{x}).
Averaging the SDE's noise produces the action
\begin{equation}
  S[\phi,\tilde\phi]
   \;=\; \int \dd t \; \Big[\,
      \tilde\phi_a\big(\Dt + \mu_a - D_a\Lap\big)\phi_a
      \;-\; \tilde\phi_a\,V_a(\{\phi\})
      \;-\; \tfrac12 \tilde\phi_a\,\Sigma_{ab}\,\tilde\phi_b
   \Big],
  \label{eq:msrjd}
\end{equation}
where $\Sigma_{ab}$ is the noise covariance (for white noise of strength $D$ this
is the local $-D\,\tilde\phi^2$ term).  Correlation and response functions of the
original SDE become correlation functions in this field theory:
\begin{equation}
  \avg{\phi_a(t)\,\phi_b(t')} \quad\text{and}\quad
  \avg{\phi_a(t)\,\tilde\phi_b(t')} = \text{(response/propagator)}.
\end{equation}
From \eqref{eq:msrjd} one reads off the two ingredients every diagrammatic
calculation needs:

\begin{itemize}
  \item the \term{bare propagators}, from the quadratic part --- a
        \emph{response} (retarded) propagator $\Gret$ that is causal
        ($\propto \theta(t-t')$), and a \emph{correlation} propagator from the
        noise vertex $\Sigma$;
  \item the \term{interaction vertices}, from the nonlinear terms
        $\tilde\phi_a V_a$ --- each monomial in the fields is a vertex with a
        definite number of legs and a prefactor.
\end{itemize}

Connected cumulants are then sums of connected Feynman diagrams built from these
propagators and vertices, organised by the number of independent loops $\ell$.
Tree diagrams ($\ell=0$) give the mean-field/linear answer; each extra loop is a
higher-order fluctuation correction.  \emph{Everything} the engine does is in
service of (i) enumerating those diagrams and (ii) evaluating their integrals.

\begin{note}
Because the response propagator is causal, the loop integrals are not the
relativistic Feynman integrals of high-energy physics; they are
\emph{time-ordered}, retarded integrals.  This is why the engine has its own
integrator (Parts IV and V) rather than reusing an off-the-shelf high-energy
loop-integral library.
\end{note}

\section{The pipeline at a glance}

A single call --- \code{compute\_cumulants} in \file{pipeline/compute.py},
wrapped for users by \code{dd.run} in \file{notebooks/daedalus.py} --- drives the
whole computation.  Internally it proceeds through seven numbered phases, which
you can watch scroll past when \code{verbose=True}:

\begin{enumerate}
  \item \textbf{expand} --- Taylor/functional-expand the symbolic action about the
        mean-field saddle and extract the interaction vertices
        (Chapter on the action and vertices).
  \item \textbf{propagator} --- build the frequency-domain propagator matrix
        $G(\omega) = K(\omega)^{-1}$ over the (physical, response) field basis
        (Chapter on the propagator).
  \item \textbf{mean\_field} --- solve the deterministic saddle-point equations
        for the steady state $\phi^{*}$ about which the expansion is taken
        (Chapter on the mean-field solver).
  \item \textbf{poles} --- find the retarded poles of $G(\omega)$ and the residue
        matrices that give the time-domain propagator (same chapter as
        propagator).
  \item \textbf{diagrams} --- enumerate the distinct diagram topologies for the
        requested $(k,\ell)$ and assign causal edge types, loading from cache
        when possible (Chapters on enumeration and typing).
  \item \textbf{classify} --- for each diagram compute its combinatorial prefactor
        and symmetry factor (Chapter on typing and symmetry).
  \item \textbf{phase\_j} --- evaluate every diagram's loop integral and sum them
        into the cumulant (Part IV for temporal models, Part V for spatial).
\end{enumerate}

\begin{figure}[t]
\centering
\begin{tikzpicture}[
   node distance=6.2mm,
   box/.style={draw,rounded corners,align=center,inner sep=4pt,
               minimum width=58mm,font=\small,fill=notebg},
   io/.style ={draw,align=center,inner sep=4pt,minimum width=58mm,
               font=\small,fill=defbg},
   arr/.style={-{Latex[length=2mm]},thick}]
  \node[io]  (txt)  {Model as text \;/\; \file{.theory.py}};
  \node[box,below=of txt] (build){\textbf{build}: parse action $\to$ symbolic model};
  \node[box,below=of build](mark) {\textbf{markovianize} (if colored noise)};
  \node[box,below=of mark] (exp)  {\textbf{expand}: vertices about the saddle};
  \node[box,below=of exp]  (prop) {\textbf{propagator} \& \textbf{mean\_field} \& \textbf{poles}};
  \node[box,below=of prop] (enum) {\textbf{diagrams}: enumerate \& type \& symmetry};
  \node[box,below=of enum] (intg) {\textbf{phase\_j}: loop integration\\(temporal \emph{or} spatial)};
  \node[box,below=of intg] (asm)  {\textbf{assemble}: cumulants / moments};
  \node[io, below=of asm]  (out)  {$\kappa_k(\tau)$, plots, saved results};
  \foreach \a/\b in {txt/build,build/mark,mark/exp,exp/prop,prop/enum,enum/intg,intg/asm,asm/out}
     \draw[arr] (\a) -- (\b);
\end{tikzpicture}
\caption{The end-to-end pipeline.  Green boxes are user-facing input/output;
blue boxes are internal phases.  The branch into a temporal or spatial integrator
happens at \textbf{phase\_j}; everything above it is shared.}
\label{fig:pipeline}
\end{figure}

Figure~\ref{fig:pipeline} shows the flow.  The crucial structural fact is that
\emph{everything above the integrator is shared between temporal and spatial
models}: the same action parser, the same vertex extraction, the same diagram
enumeration.  Only the final loop-integration phase splits into a time-domain
engine (Part~IV) and a momentum-space engine (Part~V).

\section{How the source tree is organised}

The code lives in three layers, from user-facing down to mathematical core:

\begin{description}[style=nextline,leftmargin=1.2em]
  \item[\file{notebooks/daedalus.py} --- the engine API] A single
        $\sim$1100-line convenience module that users import as \code{dd}.  It
        defines the \code{Config} dataclass, \code{load\_theory}, \code{run},
        \code{describe\_model}, and the plotting helpers.  It \emph{orchestrates};
        it contains no physics of its own (Chapter on the API).
  \item[\file{pipeline/} --- orchestration] The high-level stages:
        \file{compute.py} (the seven-phase driver), the theory-authoring API
        (\file{theory.py}, \file{theory\_compiler.py}), the propagator and
        mean-field builders (\file{\_propagator.py}, \file{\_mean\_field*.py}),
        the diagram-cache front end (\file{\_diagrams.py}), the colored-noise
        preprocessor (\file{colored\_to\_markovian.py}), and the caching
        machinery (\file{\_expand\_cache.py}, \file{\_precompute.py}).
  \item[\file{msrjd/} --- the mathematical engine] The lower-level core:
        \file{core/} (the symbolic action, vertices, convolutions),
        \file{enumeration/} and \file{diagrams/} (diagram generation, typing,
        causality, symmetry factors), and \file{integration/} which splits into
        \file{time\_domain/} (the temporal Phase-J integrator, dominated by the
        5000-line \file{final\_integral.py}) and \file{spatial/} (the
        fourteen-module momentum-space integrator).
\end{description}

Two further directories complete the picture: \file{theories/} holds the on-disk
\file{*.theory.py} model descriptions, and \file{models/} holds the
\emph{independent} simulators used to validate the diagrammatic output --- direct
Euler--Maruyama and Gillespie integrators that share no code with the pipeline
(Chapter on simulators).

\section{Caching: why a second run is instant}

Several phases are expensive but deterministic functions of a small key, so the
engine caches them on disk under \file{saved\_theories/<model>/}.  The most
important is the \term{diagram-enumeration cache}: enumerating and typing all
diagrams for a given $(k,\ell)$ is combinatorially heavy at two loops, but it
depends only on the model's topology, not on parameter values, so it is computed
once and reloaded from a Sage-object (\file{.sobj}) file thereafter.  A second
cache stores the expanded vertices.  Together they are why changing
$\mu$ and re-running is fast even though the first run of a new two-loop model is
not (Chapter on caching).

\section{Validation philosophy}

Every nontrivial result is checked against an independent Monte-Carlo
simulation.  The simulators in \file{models/} integrate the \emph{original} SDE
directly --- no diagrams, no propagators --- and estimate the same cumulant the
pipeline predicts.  Agreement between the two is the engine's ground truth; the
example notebooks (\file{notebooks/examples/}) each end with such an overlay.  As
a concrete instance, the white-noise quartic oscillator of
\eqref{eq:langevin} with $\mu=1,\varepsilon=0.02,D=1$ has an equal-time variance
the linear theory puts at $D/\mu = 1$; the two-loop pipeline corrects this to
$0.950$, and a direct simulation gives $0.946\pm0.001$ --- agreement to a fraction
of a percent.

\section{Where to go next}

The next chapter walks through a complete worked example end to end.  After that,
Part~II shows how to describe your own model, and Parts~III--V open up each stage
of the machinery in full detail.
